\chapter{Subtração}\label{ch:g2:subtraction}

Quando as unidades de cima são poucas demais para tirar delas,
desmanchamos uma dezena --- é o \emph{empréstimo}, o espelho do
vai-um. Este capítulo o explica com palitos e depois treina a
subtração armada, com a contagem para cima como atalho amigo.

\section{Por que pedimos emprestado}

\begin{example}[Desmanchar uma dezena]\label{ex:g2:subtraction:unbundle}
$42 - 17$: queremos tirar $7$ unidades de $2$ unidades --- poucas
demais! Então desmanchamos uma das $4$ dezenas em dez unidades
soltas: agora o número de cima tem $3$ dezenas e $12$ unidades.
Tirando: $12 - 7 = 5$ unidades e $3 - 1 = 2$ dezenas. Logo
$42 - 17 = 25$.
\end{example}

\begin{omfigure}
\begin{tikzpicture}[scale=0.75]
  % uma dezena amarrada
  \foreach \i in {0,...,9} \draw[omExo, very thick]
    (0.2*\i,0) -- (0.2*\i,1.7);
  \draw[omExo, thick, rounded corners=3pt]
    (-0.2,0.62) rectangle (2.0,1.05);
  \node[below, font=\small] at (0.9,-0.25) {uma dezena};
  % desmanchada em dez unidades soltas
  \draw[-Stealth, omThm, very thick] (2.6,0.85) -- (3.6,0.85);
  \foreach \i in {0,...,9} \draw[omDef, very thick]
    (4.2+0.34*\i,0) -- (4.2+0.34*\i,1.7);
  \node[below, font=\small] at (5.73,-0.25) {dez unidades soltas};
  \node[right, font=\small, align=left] at (8.0,0.85)
    {desmanche uma dezena para ter\\unidades bastantes para tirar};
\end{tikzpicture}

{\small O empréstimo é um desmanche: uma dezena vira dez unidades, e
aí a subtração pode ser feita.}
\end{omfigure}

\section{A conta armada}

\begin{method}[Subtração em coluna]\label{met:g2:subtraction:column}
\begin{enumerate}
  \item Escreva o número maior em cima, alinhado à direita;
  \item subtraia as unidades; se o algarismo de cima for pequeno
        demais, peça uma dezena emprestada (ela vira dez unidades na
        coluna atual);
  \item subtraia as dezenas e depois as centenas;
  \item \emph{confira}: o resultado mais o número que foi tirado tem
        de devolver o número de cima.
\end{enumerate}
\end{method}

\begin{example}\label{ex:g2:subtraction:worked}
Calcule $503 - 168$:
\[
\begin{array}{r}
5\,0\,3 \\
-\ 1\,6\,8 \\
\hline
3\,3\,5
\end{array}
\]
Unidades: $3 - 8$ não dá; peça emprestado passando pelas dezenas
(não há nenhuma, então peça primeiro às centenas): o número de cima
vira $4$ centenas, $9$ dezenas, $13$ unidades. Aí $13 - 8 = 5$;
$9 - 6 = 3$; $4 - 1 = 3$. Conferência: $335 + 168 = 503$.
\checkmark
\end{example}

\section{Contar para cima}

\begin{method}[Subtrair contando para cima]\label{met:g2:subtraction:countup}
Para números próximos, suba do menor até o maior e some os pulos:
\[
83 - 78: \quad 78 \to 80 \text{ são } 2, \ 80 \to 83 \text{ são } 3,
\text{ logo } 83 - 78 = 5 .
\]
Contar para cima é mais rápido que pedir emprestado quando os dois
números estão perto um do outro --- e é exatamente assim que o
lojista dá o troco.
\end{method}

\begin{omfigure}
\begin{tikzpicture}[scale=0.85]
  \draw[-Stealth] (76.6,0) -- (84.4,0);
  \foreach \i in {77,...,84} \draw (\i,-0.1) -- (\i,0.1);
  \foreach \i in {78,80,83}
    \node[below, font=\footnotesize] at (\i,-0.12) {\i};
  \draw[-Stealth, omDef, thick] (78,0.3) .. controls (79,0.9) ..
    (80,0.3);
  \node[omDef, font=\small] at (79,1.0) {$+2$};
  \draw[-Stealth, omThm, thick] (80,0.3) .. controls (81.5,1.05) ..
    (83,0.3);
  \node[omThm, font=\small] at (81.5,1.1) {$+3$};
\end{tikzpicture}

{\small $83 - 78$ contando para cima: dois pulos, $2 + 3 = 5$.}
\end{omfigure}

\begin{example}[Dar o troco]\label{ex:g2:subtraction:change}
Um brinquedo custa $34$; você paga com $50$. Troco: conte para cima,
$34 \to 40$ são $6$, $40 \to 50$ são $10$, então o troco é $16$. A
subtração $50 - 34 = 16$, feita do jeito do lojista.
\end{example}

\section{Exercícios}

\begin{exercise}[$\star$]\label{exo:g2:subtraction:1}
Calcule (sem precisar de empréstimo): $58 - 23$; \quad $176 - 45$;
\quad $389 - 164$.
\end{exercise}

\begin{exercise}[$\star$]\label{exo:g2:subtraction:2}
Calcule com empréstimo: $52 - 17$; \quad $73 - 28$; \quad
$60 - 24$.
\end{exercise}

\begin{exercise}[$\star$]\label{exo:g2:subtraction:3}
Arme a conta e confira cada uma: $324 - 167$; \quad
$500 - 236$; \quad $703 - 158$.
\end{exercise}

\begin{exercise}[$\star$]\label{exo:g2:subtraction:4}
Calcule contando para cima: $62 - 58$; \quad $91 - 85$; \quad
$100 - 96$.
\end{exercise}

\begin{exercise}[$\star$]\label{exo:g2:subtraction:5}
Calcule o troco (contando para cima): pague $50$ por um livro de
$42$; pague $20$ por um brinquedo de $13$; pague $100$ por uma bolsa
de $75$.
\end{exercise}

\begin{exercise}[$\star$]\label{exo:g2:subtraction:6}
``Uma árvore tem $156$ cm de altura. Ela cresceu $28$ cm este ano.
Que altura tinha no ano passado?'' Escreva a subtração e responda.
\end{exercise}

\begin{exercise}[$\star$]\label{exo:g2:subtraction:7}
Confira com a adição gêmea: $85 - 39 = 46$; está certo?
$214 - 68 = 156$; está certo?
\end{exercise}

\begin{exercise}[$\star$]\label{exo:g2:subtraction:8}
Para cada história, escolha $+$ ou $-$: ``$47$ passarinhos, $19$
voam para longe''; ``$47$ maçãs, mais $19$ colhidas''; ``Samuel tem
$47$, Lia tem $19$: quantos Samuel tem a mais?''.
\end{exercise}

\begin{exercise}[$\star$]\label{exo:g2:subtraction:9}
Preencha o buraco: $63 - \;?\; = 45$. (Conte para cima, de $45$ até
$63$.)
\end{exercise}

\begin{exercise}[$\star\star$]\label{exo:g2:subtraction:10}
Uma brincadeira: pense num número, some $30$ e depois tire $18$. O
resultado é $58$. Qual era o número? (Volte de trás para frente!)
\end{exercise}
